\documentclass[11pt,a4paper]{article}
\usepackage[utf8]{inputenc}
\usepackage[T1]{fontenc}
\usepackage[spanish]{babel}
\usepackage{geometry}
\geometry{letterpaper, margin=1in}
\usepackage{setspace}
\setstretch{1.2}
\usepackage{graphicx}
\usepackage{booktabs}
\usepackage{float}
\usepackage{hyperref}
\hypersetup{colorlinks=true, linkcolor=blue, urlcolor=blue, citecolor=blue}
\usepackage[table]{xcolor}
\usepackage{longtable}
\usepackage{enumitem}
\usepackage{fancyhdr}
\pagestyle{fancy}
\fancyhf{}
\fancyhead[L]{SRS v1.0.0 --- Sistema de Gestión de Pre-Sustentaciones UTEQ}
\fancyhead[R]{\thepage}
\renewcommand{\headrulewidth}{0.4pt}

\title{}
\date{}

\begin{document}

% ============ CARATULA ============
\begin{titlepage}
\centering
\vspace*{1cm}
{\LARGE \textbf{UNIVERSIDAD TÉCNICA ESTATAL DE QUEVEDO}}\\[0.3cm]
{\large Facultad de Ciencias de la Computación}\\[0.2cm]
{\large Carrera de Ingeniería de Software (Rediseño)}\\[0.8cm]

{\large \textbf{ASIGNATURA: Aplicaciones Web}}\\[0.3cm]
{\large Quinto Nivel --- Período Académico 2026-2027 PPA}\\[1.2cm]

{\Large \textbf{ESPECIFICACIÓN DE REQUISITOS DE SOFTWARE (SRS)}}\\[0.3cm]
{\large Estándar ISO/IEC/IEEE 29148:2018}\\[0.6cm]
{\LARGE \textbf{Sistema de Gestión de Pre-Sustentaciones UTEQ}}\\[0.6cm]
{\Large \textbf{Versión Final v1.0.0}}\\[1.2cm]

\vfill
{\large \textbf{EQUIPO DE DESARROLLO:}}\\[0.3cm]
{\large Alava Alvarado, Jean Pierre \quad (\href{https://orcid.org/0009-0001-2878-2919}{ORCID: 0009-0001-2878-2919})}\\[0.1cm]
{\large Moncayo Loor, Xavier Alejandro}\\[0.1cm]
{\large Zamora Arias, Carla Esthefania}\\[0.1cm]
{\large Barreto Rosado, Heider Dominick}\\[1cm]

{\large \textbf{DOCENTE-DIRECTOR:}}\\[0.3cm]
{\large Ing. Guerrero Ulloa Gleiston Cíceron, Mg.}\\[1cm]

{\large \textbf{FECHA:} 17 de agosto de 2026}\\[0.5cm]
\end{titlepage}

\tableofcontents
\newpage

% ============ HISTORIAL DE VERSIONES ============
\section{Historial de versiones}

\begin{table}[H]
\centering
\small
\begin{tabular}{p{1.5cm} p{1.7cm} p{3.2cm} p{7.5cm}}
\toprule
\textbf{Versión} & \textbf{Fecha} & \textbf{Autor} & \textbf{Cambios} \\
\midrule
0.9.0-rc & 2026-07-30 & Equipo de desarrollo & Versión inicial (Entrega 3): 5 HU formalizadas de 12 requisitos ya referenciados en la matriz de trazabilidad; sin casos de uso. \\
\addlinespace
1.0.0 & 2026-08-17 & Equipo de desarrollo, asistido por Claude/Anthropic (ver \href{../etica/ai-disclosure.md}{declaración de uso de IA}) & Versión final (Fase 7): 12 HU completas en formato Connextra+INVEST+Gherkin, 12 casos de uso Cockburn, checklist INCOSE completo, matriz de trazabilidad corregida (5 citas de evidencia falsas identificadas y corregidas), tasa de estabilidad calculada. \\
\bottomrule
\end{tabular}
\caption{Historial de versiones del SRS. Ver bitácora completa en \texttt{CHANGELOG-REQ.md}.}
\end{table}

\textbf{Nota sobre el criterio D0R:} este documento reemplaza a \texttt{docs/srs/SRS.md} (ahora archivado
en \texttt{docs/requisitos/historico/SRS-v0.9.0-rc.md}) como la especificación de requisitos vigente. La
página de aprobación (Sección~9) está lista para firma física o electrónica del docente-director; a la
fecha de generación de este PDF, \textbf{esa firma sigue pendiente} — ver la nota explícita en esa
sección.

% ============ 1. INTRODUCCION ============
\section{Introducción}

\subsection{Propósito}
El presente documento define los requisitos funcionales y no funcionales para el Sistema de Gestión de
Pre-Sustentaciones UTEQ, elaborado bajo el estándar internacional ISO/IEC/IEEE 29148:2018. Sirve como
contrato técnico entre la coordinación académica, docentes, estudiantes y el equipo desarrollador.

\subsection{Alcance del sistema}
El sistema automatiza el flujo completo de pre-sustentación de anteproyectos de titulación: registro de
solicitudes, carga de anteproyectos con verificación de integridad SHA-256, asignación de jurados,
programación de cronogramas, evaluación por rúbrica, y emisión y firma digital multi-actor de actas.

\subsection{Definiciones, acrónimos y abreviaturas}
\begin{itemize}[nosep]
    \item \textbf{SRS:} Software Requirements Specification (ISO/IEC/IEEE 29148:2018).
    \item \textbf{UTEQ:} Universidad Técnica Estatal de Quevedo.
    \item \textbf{HU:} Historia de Usuario (formato Connextra).
    \item \textbf{CU:} Caso de Uso (plantilla Cockburn).
    \item \textbf{RF / RNF:} Requisito Funcional / No Funcional.
    \item \textbf{JWT:} JSON Web Token.
    \item \textbf{SP:} Stored Procedure (procedimiento almacenado PL/pgSQL).
    \item \textbf{MoSCoW:} técnica de priorización Must/Should/Could/Won't.
    \item \textbf{INVEST:} Independent, Negotiable, Valuable, Estimable, Small, Testable.
\end{itemize}

% ============ 2. DESCRIPCION GENERAL ============
\section{Descripción general}

\subsection{Perspectiva del producto}
Solución web de arquitectura N-capas: Frontend Angular 21 (SPA), Backend Spring Boot 3.2.1, PostgreSQL
15, caché Redis 7, desplegado con Docker Compose y nginx como proxy inverso.

\subsection{Características de los usuarios}
\begin{table}[H]
\centering
\begin{tabular}{ll p{7cm}}
\toprule
\textbf{Rol} & \textbf{Descripción} & \textbf{Permisos clave} \\
\midrule
Estudiante & Postulante a titulación & Registrar solicitud, subir anteproyecto, consultar estado y acta \\
Docente/Jurado & Evaluador asignado & Revisar anteproyecto, calificar rúbrica \\
Presidente del Tribunal & Líder de la defensa & Generar y firmar acta \\
Coordinador & Gestor académico & Asignar jurados, programar cronograma, salas, reportes \\
Administrador & Admin. del sistema & Gestionar usuarios, catálogos, seguridad \\
\bottomrule
\end{tabular}
\end{table}

% ============ 3. REQUISITOS ESPECIFICOS ============
\section{Requisitos específicos}

\subsection{Requisitos funcionales (resumen)}

Los 12 requisitos funcionales se documentan en detalle, con formato Connextra, checklist INVEST y
criterios de aceptación Gherkin, en \texttt{docs/requisitos/historias/HU-01.md} a \texttt{HU-12.md}. Los
12 casos de uso correspondientes (plantilla Cockburn, 4 niveles de meta, diagrama de secuencia) están en
\texttt{docs/requisitos/casos-de-uso/CU-01.md} a \texttt{CU-12.md}. Esta sección resume su alcance y
estado de verificación real.

\begin{longtable}{p{1.2cm} p{3.8cm} p{1.6cm} p{2.2cm} p{3cm}}
\toprule
\textbf{RF} & \textbf{Título} & \textbf{Prioridad} & \textbf{Módulo} & \textbf{Verificación} \\
\midrule
\endhead
RF-01 & Autenticación segura & Must & Auth & \textcolor{green!50!black}{Verificado (test real)} \\
RF-02 & Registro de solicitud & Must & Solicitudes & \textcolor{green!50!black}{Verificado (test real)} \\
RF-03 & Asignación de jurados & Must & Jurados & \textcolor{green!50!black}{Verificado (test real)} \\
RF-04 & Programación de cronograma & Must & Cronograma & \textcolor{green!50!black}{Verificado (test real)} \\
RF-05 & Evaluación por rúbrica & Must & Evaluaciones & \textcolor{green!50!black}{Verificado (test real)} \\
RF-06 & Generación de actas & Must & Actas & \textcolor{green!50!black}{Verificado (test real)} \\
RF-07 & Firma digital de actas & Should & Actas/Firmas & \textcolor{green!50!black}{Verificado (test real)} \\
RF-08 & Notificaciones & Could & Notificaciones & \textcolor{red}{No verificado} \\
RF-09 & Reportes de defensas & Could & Reportes & \textcolor{red}{No verificado} \\
RF-10 & Gestión de salas & Should & Salas & \textcolor{red}{No verificado} \\
RF-11 & Gestión de usuarios & Must & Usuarios & \textcolor{green!50!black}{Verificado (test real)} \\
RF-12 & Carga de anteproyecto & Must & Anteproyectos & \textcolor{green!50!black}{Verificado (test real)} \\
\bottomrule
\caption{Resumen de los 12 requisitos funcionales. Fuente de verdad: \texttt{docs/trazabilidad/matriz.csv} v1.0.0.}
\end{longtable}

\textbf{Estado real del criterio D0R (100\% de los Must verificados):} de los 8 requisitos con
prioridad \textbf{Must}, \textbf{los 8 tienen prueba automatizada real (100\%, actualizado 2026-08-29)}.
Esta cifra se calculó ejecutando \texttt{scripts/validate-traceability.sh} contra el repositorio real, no
se estimó. La cifra subió desde el 5/8 (62.5\%) original de la Fase 6 tras añadir
\texttt{SolicitudServiceImplTest.java} (RF-02), \texttt{CronogramaServiceImplTest.java} (RF-04) y,
finalmente, \texttt{ActaServiceImplTest.java} (RF-06, el último que faltaba). El 100\% aplica a los
requisitos \textbf{Must} exigidos por D0R, no a los 12 requisitos totales del sistema: RF-08, RF-09 y
RF-10 (prioridad Could/Should) siguen sin prueba automatizada dedicada (9/12 = 75\% general) --- se
declara esta distinción explícitamente en vez de reportar un 100\% general que la matriz no respalda.

\subsection{Requisitos no funcionales}

\begin{table}[H]
\centering
\begin{tabular}{ll p{8cm}}
\toprule
\textbf{RNF} & \textbf{Categoría} & \textbf{Descripción y estado real} \\
\midrule
RNF-01 & Rendimiento & Peticiones REST $<$ 500ms bajo 50 usuarios concurrentes. \textbf{Verificado}: p95 real de 7.4--80.9ms en 5 corridas k6 (\texttt{k6/README.md}). \\
RNF-02 & Seguridad & Protección OWASP Top 10. \textbf{Verificado parcialmente}: 0 hallazgos de SQL dinámico (find-sec-bugs), 0 FAIL y 0 High en OWASP ZAP (corregido 2026-08-29), 14 dependencias vulnerables sin corregir en tooling de build (\texttt{npm audit}, bajó de 41 tras corregir \texttt{@angular/core}). \\
RNF-03 & Usabilidad & Puntaje SUS $>$ 75/100. \textbf{No verificado}: instrumento listo, sin aplicar a usuarios reales (\texttt{docs/mediciones/sus/SUS-RESULTS.md}). \\
RNF-04 & Mantenibilidad & Cobertura JaCoCo $\geq$ 60\%. \textbf{No cumplido}: 38.88\% de líneas real (2026-08-31, 109 tests), en progreso desde 2.83\% (\texttt{docs/mediciones/jacoco/COVERAGE.md}). \\
\bottomrule
\end{tabular}
\end{table}

% ============ 4. TRAZABILIDAD ============
\section{Trazabilidad y verificación (Capítulo dedicado, no fusionado con diseño)}

La matriz de trazabilidad completa (\texttt{docs/trazabilidad/matriz.csv} v1.0.0) conecta cada RF con su
HU, módulo backend, endpoint REST, prioridad MoSCoW, prueba automatizada real (o su ausencia declarada) y
evidencia empírica verificada. Se corrigieron en esta versión 5 citas de archivos de test inexistentes y
3 citas de evidencia empírica incorrectas (\textit{non sequiturs} como usar el puntaje SUS general como
evidencia de una funcionalidad CRUD específica) heredadas de la matriz v0.9.0-rc --- ver el detalle
completo en \texttt{CHANGELOG-REQ.md}.

% ============ 5. ANEXO INCOSE ============
\section{Anexo: Checklist INCOSE v4}

Se aplicaron las 9 características de calidad individuales de INCOSE (Necesario, Apropiado, No ambiguo,
Completo, Singular, Factible, Verificable, Correcto, Conforme) a los 12 requisitos, más las 6
características de calidad del conjunto (Completo, Consistente, Factible, Comprensible, Verificable como
conjunto, Acotado). El checklist completo, con la matriz de 12$\times$9 y el análisis de conjunto, está
en \texttt{docs/checklists/INCOSE-REQUIREMENTS.md} y se anexa por referencia a este documento en vez de
duplicarse aquí verbatim, para evitar que las dos copias diverjan con el tiempo.

\textbf{Conclusión del anexo:} la brecha más consistente es la característica ``Conforme'' (formato
Historia de Usuario en vez de lenguaje normativo ``shall'' --- decisión de estilo consciente) y
``Correcto''/``Verificable como conjunto'' para los 7 requisitos sin prueba automatizada real.

% ============ 6. ELICITACION ============
\section{Anexo: Técnicas de elicitación aplicadas}

Resumen (detalle completo en \texttt{docs/requisitos/elicitacion/README.md}): de las 5 técnicas
esperadas, \textbf{2 tienen evidencia real y verificable} en el repositorio (análisis de documentos ---
7 observaciones docentes trazadas a commits reales en \texttt{OBSERVACIONES.md}, y evaluación cruzada de
pares en \texttt{AUTOEVALUACION-PRESUS.md}), \textbf{1 es parcial} (prototipado iterativo vía historial
de Git, sin sesiones formales documentadas), y \textbf{2 no tienen evidencia documentada} (entrevistas,
observación directa, workshops) --- se declara esto explícitamente en vez de fabricar transcripciones o
actas que no existieron.

% ============ 7. PAGINA DE APROBACION ============
\section{Página de aprobación}

\begin{center}
\begin{tabular}{p{7cm} p{7cm}}
\toprule
\textbf{Elaborado por (Equipo de desarrollo)} & \textbf{Revisado y aprobado por (Docente-Director)} \\
\midrule
& \\[1.5cm]
\rule{6cm}{0.4pt} & \rule{6cm}{0.4pt} \\
Alava Alvarado, Jean Pierre & Ing. Guerrero Ulloa Gleiston Cíceron, Mg. \\
\vspace{0.3cm} & \\
\rule{6cm}{0.4pt} & Fecha de aprobación: \rule{4cm}{0.4pt} \\
Moncayo Loor, Xavier Alejandro & \\
\vspace{0.3cm} & \\
\rule{6cm}{0.4pt} & \\
Zamora Arias, Carla Esthefania & \\
\vspace{0.3cm} & \\
\rule{6cm}{0.4pt} & \\
Barreto Rosado, Heider Dominick & \\
\bottomrule
\end{tabular}
\end{center}

\vspace{0.5cm}
\noindent\fbox{\parbox{\textwidth}{\textbf{Estado real de la firma (2026-08-17):} este documento se
generó de forma automatizada como parte del cierre de la ingeniería de requisitos (Fase 7). La firma del
docente-director \textbf{no se ha obtenido todavía} --- requiere la revisión y aprobación real del Ing.
Guerrero Ulloa por parte del equipo humano, fuera del alcance de lo que se puede generar
automáticamente. No se fabrica una firma ni se marca esta sección como completa hasta que ocurra
realmente, siguiendo la misma política de integridad aplicada al resto de este repositorio (ver
\href{../etica/ai-disclosure.md}{declaración de uso de IA}).}}

\end{document}
